%%%%%%%%%%%%%%%%%%%%%%%%%%%%%%%%%%%%%%%%%%%%%%%%%%%%%%%%%%%%%%
%%%%%%%%%%%%%%%%% MA-XXXX Análisis Real II %%%%%%%%%%%%%%%%%%%
%%%%%%%%%%%%%%%%%%%%%%%%%%%%%%%%%%%%%%%%%%%%%%%%%%%%%%%%%%%%%%
\newpage
\subsection*{MA-0725 Análisis Real II}
\addcontentsline{toc}{subsection}{MA-0725 Análisis Real II}

\hypertarget{MA0725}{}
\begin{refsection}
\begin{table}[H]
\centering{
\begin{tabular}{|ll|}
\hline
%\multicolumn{2}{|c|}{\textbf{Nivel de virtualidad del curso:} Virtual}\\ \hline
\textbf{Año:} IV  & \textbf{Requisitos:} Por determinar \\ 
\textbf{Ciclo:} VII  & \textbf{Correquisitos:} Ninguno \\ 
\textbf{Tipo de curso:} Teórico & \textbf{Horas presenciales por semana:}   \\ 
\textbf{Créditos:} 4 & \textbf{Horas de trabajo independiente por semana:}  \\ \hline
\end{tabular}
}
\end{table}



\subsubsection*{Descripción del curso}

Este es un curso del cuarto año de la carrera Bach. y Lic. en Matemática en el que se continua el desarrollo de las bases del Análisis Real, con un enfoque hacia el estudio de los espacios clásicos de Banach y una introducción al Análisis de Fourier.

Este curso se enfocará en fortalecer el rigor en el tratamiento y la comunicación de argumentos matemáticos, así como en la presentación de demostraciones de manera clara y precisa.

La primera parte del curso cubre la Teoría de Diferenciación de Lebesgue, seguida de la noción de continuidad absoluta y el Teorema Fundamental del Cálculo. La segunda parte tratará de los espacios $L^p$ y sus propiedades. Luego, se tratarán algunos elementos básicos de la teoría de espacios de Banach y Hilbert. Por último, se dará una introducción a las series y a la transformada de Fourier.

Para el buen desempeño en este curso es necesario tener conocimientos previos de Algebra Lineal, Topología General, así como de la Teoría de la Medida e Integración de Lebesgue. 



\subsubsection*{Objetivos}

\textbf{General:} Aplicar los conceptos y resultados intermedios del Análisis Real para la resolución de problemas. \\

\textbf{Específicos:} Al finalizar el curso se espera que la persona estudiante sea capaz de:

\begin{enumerate}
\item Resolver problemas matemáticos y de disciplinas afines utilizando los conceptos y resultados básicos del Análisis de Fourier.
\item Aplicar los conceptos b\'asicos de diferenciación y continudad absoluta, en el planteamiento y resoluci\'on de problemas matem\'aticos.
\item Demostrar resultados sobre propiedades de los espacios clásicos de Banach.
\item Indagar en la literatura sobre temas que permitan complementar su formación.
\item Comunicar ideas matemáticas de manera clara y efectiva en forma oral y escrita.
\end{enumerate}

\subsubsection*{Contenidos}

\begin{enumerate}
\item \textbf{Diferenciación (3 semanas):} 
\begin{enumerate}
    \item Lemas de cubrimiento de Vitalli.
    \item Diferenciación de funciones monótonas. Números de Dini.
    \item Convergencia en promedio, funciones localmente integrables.
    \item La Función Maximal de Hardy-Littlewood y el Teorema Maximal.
    \item Teorema de Diferenciación de Lebesgue, conjunto de Lebesgue.
    \item Convexidad.
\end{enumerate}

\item \textbf{Continuidad absoluta y el Teorema Fundamental del Cálculo (2 semanas):} 

\begin{enumerate}
    \item Continuidad absoluta: definición, propiedades, diferenciabilidad.
    \item Teorema de Banach-Zaretsky.
    \item Teorema Fundamental del Cálculo y aplicaciones.
\end{enumerate}



\item \textbf{Espacios de Banach (2 semanas):} 
\begin{enumerate}
    \item Normas: Definición y propiedades, espacios normados, métrica inducida.
    \item Completitud y Espacios de Banach.
    \item Convergencia de funciones, convergencia uniforme.
    \item Equivalencia de normas.
    \item Dimensión y espacios de dimensión finita. Compacidad local.
    \item Operadores lineales.
    \item Series de funciones y convergencia absoluta.
\end{enumerate}

\item \textbf{Espacios de Hilbert (2 semanas):} 
\begin{enumerate}
    \item Espacios de producto interno: Definición y propiedades.
    \item Espacios de Hilbert.
    \item Ortogonalidad: complementos, proyecciones, sucesiones completas.
    \item Sucesiones y bases ortonormales: familias ortonormales, proyecciones, existencia de bases ortonormales.
\end{enumerate}



\item \textbf{Introducción al Análisis de Fourier (4 semanas):} 
\begin{enumerate}
    \item Series de Fourier: definición, propiedades, métodos de suma, núcleos.
    \item Convergencia de series de Fourier: en $L^2$, puntual y uniforme.
    \item Transformada de Fourier: Definición, propiedades, convergencia.
\end{enumerate}



\end{enumerate}

\subsubsection*{Metodología}

\noindent \textbf{Lineamientos metodológicos:} La persona docente a cargo de este curso debe respetar las siguientes pautas:
\begin{enumerate}
    \item Fomentar el desarrollo de intuición sobre los conceptos estudiados en el curso. 
    \item Requerir de parte del estudiantado, una comunicación clara y rigurosa de argumentos y pruebas matemáticas de manera oral y escrita.
    \item Cuando la naturaleza del contenido lo permita, se deben utilizar representaciones visuales que apoyen la comprensión de los temas.
    \item Fomentar una visión integral de las matemáticas y de cómo se enlazan las diferentes áreas a través del análisis real.
    \item Para fomentar el desarrollo de las habilidades investigativas en el estudiantado, se deben implementar proyectos de investigación y si es posible la participación en coloquios o seminarios de investigación.
\end{enumerate}

\textbf{Sugerencias metodológicas:} Adicionalmente, se tienen las siguientes recomendaciones para la persona docente a cargo de este curso:
\begin{enumerate}
    \item Dado que algunos de los temas permiten cubrirse con distintos grados de profundidad en su presentación, se le recomienda a la persona docente considerar esto en su planificación del curso para no sobrecargar al estudiantado. Entre estos temas están: 
    \begin{itemize}
        \item La prueba del Teorema de Banach-Zaretski.
        \item Aproximaciones de la identidad.
        \item Convergencia puntual de series de Fourier.
    \end{itemize}
    \item Utilizar el recurso tecnológico para reforzar distintos conceptos estudiados en el curso por medio de la visualización, cálculo y experimentación.
    \item Aprovechar momentos adecuados del curso para motivar el estudio por medio de reseñas acerca del desarrollo histórico del análisis real y su importancia en aplicaciones dentro y fuera de la matemática. Por ejemplo se puede consultar \cite{Fol16} y \cite{Bre08}.
    \item En la evaluación, se sugiere utilizar estrategias que promuevan el trabajo constante de parte de las personas estudiantes a lo largo del ciclo lectivo. Por ejemplo, esto se puede hacer por medio de listas de ejercicios recomendados o proyectos.
    \item Dado que hoy en día la mayoría del trabajo en matemática, tanto a nivel académico como profesional, es de naturaleza colaborativa, se sugiere a la persona docente implementar estrategias que promuevan el trabajo en equipo.
\end{enumerate}



%\nocite{Apo76}
\nocite{Lie01}
\nocite{RF10}
%\nocite{Roy88}
\nocite{WZ15}
\nocite{SS05}
\nocite{Hei19}
\nocite{BBT08}
\nocite{Bre08}
\nocite{Fol99}
\nocite{Hei18}
\nocite{Duo01}
\nocite{Kat04}
\nocite{Fol92}
\nocite{SS103}
\nocite{DM72}




\printbibliography[heading=subbibliography,section=\the\value{refsection}]
\end{refsection}
